\section{Introduction}
\label{sec:introduction}

\subsection{Why This Problem Matters}
A token holder can exit a position in seconds. An operator who has bought, installed, and calibrated site-specific hardware cannot. In DePIN, a rooftop GNSS station, wireless device, or other physical asset is not rebalanced with the same ease as software-mediated capital. Once infrastructure is deployed, exit and re-entry become asymmetric: participation can deteriorate quickly when profitability disappears, while restoring capacity is slower because it requires renewed capital, logistics, and operating effort \cite{Ballandies2023, Lin2024, Arthur1989}.

This speed mismatch is what makes DePIN analytically different from software-only crypto systems. DePIN networks coordinate real-world infrastructure through token-based incentives, but the token is not merely a peripheral exchange layer. It is the coordination mechanism that determines whether infrastructure is deployed, retained, and kept serviceable under changing market conditions \cite{Messari2024, FrontiersDePIN2025}. When incentives weaken, providers do not adjust only by moving tokens; they may defer maintenance, reduce operational commitment, or eventually power down hardware altogether.

Much of the public discussion around DePIN still emphasizes adoption, network growth, and token appreciation. Growth can be visible long before underlying robustness is tested. Those themes matter, but they do not answer the harder operational question: what happens when demand weakens, liquidity tightens, provider costs rise, or adjacent networks offer better yields for similar hardware \cite{Ho2022, Chiu2024}? For hardware-dependent systems, robustness is not mainly an equilibrium story. It is a stress story.

\subsection{Problem Setting}
We therefore treat DePIN tokenomics as a cyber-physical coordination problem rather than as a purely financial mechanism. Under adverse conditions, the incentive system must continue to support three coupled functions:
\begin{enumerate}
    \item \textbf{Provider retention}, so participation does not erode faster than the network can absorb;
    \item \textbf{Service continuity}, so usable capacity and output remain available; and
    \item \textbf{Incentive--usage alignment}, so reward outflows do not remain detached from usage-linked value capture for too long.
\end{enumerate}

Those functions do not necessarily fail at the same time. A demand contraction can weaken usage-linked sinks before participation visibly falls. A liquidity shock can compress provider margins before any service metric moves. Competitive yield pressure can pull supply away from a network even when internal demand has not yet collapsed. For that reason, DePIN stress evaluation is less about assigning static mechanism labels and more about tracing transmission paths: where deviations appear first, how they propagate, and when they become operationally meaningful \cite{Morris2019, Braakman2022}.

\subsection{Research Question and Contribution}
\label{subsec:research_question}
The preceding argument narrows the problem in a specific way. If DePIN robustness depends on how stress propagates through provider retention, service continuity, and incentive--usage alignment, then our central task is not simply to label tokenomic designs as stable or unstable. We instead ask how different mechanisms transmit stress, where those pressures register first, and how robustness can be evaluated before breakdown becomes clearly visible.

We organize the analysis around one primary research question and one subordinate evaluative question:

\begin{quote}
How do different DePIN tokenomic architectures, specifically Burn-and-Mint Equilibrium (BME)-oriented and capped-supply designs, transmit physical and economic stress through provider participation, token markets, and service utilization?
\end{quote}

\begin{quote}
How can robustness be evaluated before failure becomes operationally visible?
\end{quote}

Our answer is comparative rather than predictive. We do not aim to forecast token prices or identify a universal winner among mechanism classes. Instead, we examine where stress concentrates first, how deviations propagate across provider retention, service continuity, and incentive--usage alignment, and how those sequences can be evaluated before breakdown becomes clearly visible under standardized stress conditions.

This framing yields three contributions. First, we shift the evaluation of DePIN tokenomics away from growth-phase narrative and toward stress transmission under physical constraints \cite{Ballandies2023, Lin2024}. Second, we use the Onocoy network as an empirical anchor for translating documented mechanism rules into a tractable evaluative setup. Third, we develop and apply DTSE, the DePIN Tokenomic Stress Evaluator, as a reproducible comparative framework for evaluating how alternative mechanism profiles react to matched stress channels under explicit assumptions \cite{Morris2019, Braakman2022}.

\subsection{Scope and Thesis Roadmap}
Our scope is deliberately bounded. We do not claim to measure live-network behavior directly, and we do not treat simulation outputs as forecasts. Instead, we combine documented mechanism facts, empirical stress context, explicit model assumptions, and DTSE outputs to support comparative statements about sensitivity, failure signatures, and robustness within a defined evaluative boundary.

The remainder of the thesis follows that same sequence. We first establish why DePIN is analytically distinct from software-only crypto systems, then develop the comparative tokenomic vocabulary needed for stress evaluation. We next ground the analysis in the Onocoy case, clarify what can and cannot be observed empirically, formalize DTSE, and finally use that framework to evaluate stress behavior and interpret its implications.
